%-------------------------------------------------------------------------------
%	ABSTRACT PAGE
%-------------------------------------------------------------------------------

\addchaptertocentry{Хураангуй} % Хураангуйг гарчигт нэмэх

\begin{center}
{\scshape\Large \univname\par} % Их сургуулийн нэр
{\scshape\large \facname\par}\vspace{0.5cm}
{\huge\textbf{{Хураангуй}} \par}
\bigskip
{\Large{\ttitle} \par} % Тезисийн нэр
\bigskip

{\normalsize \shortname \par} % Зохиогчийн нэр
\mailaddress
\end{center}

\bigskip

\textit{\textbf{Түлхүүр үгс:} Flutter, Firebase, Node.js, Docker, AWS EC2, GitHub Actions, AI чатбот, цаг захиалга, спорт заал, үүлэн орчин, Даланзадгад}

\bigskip

Энэхүү бакалаврын төгсөлтийн ажлын зорилго нь Өмнөговь аймгийн Даланзадгад хотын таван спорт заалд зориулсан үүлэн орчинд суурилсан цаг захиалгын мобайл аппликейшн хөгжүүлэх явдал юм.

Систем нь \textbf{Flutter} (Android/iOS), \textbf{Node.js/Express} бэкэнд, \textbf{Firebase} (Auth, Firestore, Cloud Functions), \textbf{Docker/AWS EC2} дэд бүтэц болон \textbf{AI чатбот} (Claude Haiku, Groq Llama) гэсэн таван бүрэлдэхүүнтэй гурван давхаргат архитектурт суурилна. Бэкэнд нь \textbf{AWS EC2} (ap-southeast-1) дээр Docker контейнераар ажиллаж, \textbf{GitHub Actions} CI/CD хоолойгоор автоматаар байршуулагддаг.

Үндсэн функцүүд: email/нууц үгийн нэвтрэлт (Nodemailer 6 оронтой кодын баталгаажуулалттай), цаг захиалах болон цуцлах (DBK-xxxxx кодтой), байгууллагуудын тогтмол захиалга (\code{fixed\_bookings}), 1 цагийн өмнөх локал мэдэгдэл, AI чатбот.

Функциональ туршилтаар 20 тест тохиолдол бүгд амжилттай давж, хэрэглэхэд хялбар байдлын нийт оноо \textbf{4.3/5.0} болсон.

\bigskip

\textit{\textbf{Keywords:} Flutter, Firebase, Node.js, Docker, AWS EC2, GitHub Actions, AI chatbot, booking system, sports hall, cloud, Dalanzadgad}

\bigskip

This bachelor's thesis develops a cloud-based three-tier sports hall booking mobile application for Dalanzadgad city, Ömnögovi Province. The system comprises a Flutter mobile client (Android/iOS), a Node.js/Express backend on AWS EC2 via Docker, Firebase services (Authentication, Firestore, Cloud Functions), and dual AI chatbot (Anthropic Claude Haiku 4.5 / Groq Llama 3.1-8b). Automated deployment is handled by GitHub Actions CI/CD. All 20 test cases passed with a usability score of 4.3/5.0.
